% Part: counterfactuals
% Chapter: minimal-change-semantics
% Section: introduction

\documentclass[../../../include/open-logic-section]{subfiles}

\begin{document}

\olfileid{cnt}{min}{int}

\olsection{پېژندنه}

سټالنکر او لويس د واقعيت خلاف شرطونو، لکه «که د اورلګيت لرګے
وهل شوے وے، نو بل شوے به وے»، شننې وړاندې کړې۔ د هغو شننې د
داسې جملو د صدق-شرطونو د سم پوهېدو وړانديزونه وو۔ د دواړو
وړانديزونو تر شا نظر دا دے: د دې د ارزولو لپاره چې د واقعيت
خلاف شرطي رښتينې ده که نۀ، بايد هغه ممکنې نړۍ په پام کښې
ونيسو چې د مقدم د رښتيني کولو لپاره د واقعي نړۍ له حالت څخه
تر ټولو لږې بېلې دي۔ که پايله په دغو ممکنو نړۍو کښې رښتينې
وي، نو د واقعيت خلاف شرطي رښتينې ده۔ د بېلګې په توګه، فرض کړئ
چې زما په لاس کښې د اورلګيت يو لرګے او د اورلګيتونو يوه کتابچه
ده۔ په واقعي نړۍ کښې زه يوازې ورته ګورم او سوچ کوم چې که
لرګے ووهم نو څه به وشي۔ له واقعي نړۍ څخه هغه لږ تر لږه بدلون
چې پکښې لرګے وهم، دا دے چې د عمل پرېکړه وکړم او لرګے ووهم۔
بدلون په دې معنا لږ تر لږه دے چې نور هېڅ نۀ بدلېږي: زه ورسره
په هوا کښې ټوپ نۀ وهم، د لرګي وهل زما وېښتۀ هم نۀ بلوي، زه
ناڅاپه په ګوتو کښې ټول ځواک نۀ بايلم، په هماغه وخت د اوبو د
لوبې د سوپر سوکر ټوپک په کمين کښې پر ما اوبۀ نۀ اچول کېږي، او داسې نور۔
په دغه بديل امکان کښې د اورلګيت لرګے بلېږي۔ نو دا رښتينې ده
چې که زه لرګے ووهم، نو بل به شي۔

د دې لومړني تصور شننه د واقعيت خلاف شرطونو د منطقونو له رسمي
معنٰی پوهنې سره يو ځاے کېدے شي۔ لويس د واقعيت خلاف شرطي لپاره
«$\boxright$» نښه معرفي کړه، او سټالنکر «$>$» نښه کاروله۔
مونږ به~$\cif$ کاروو او د ويناوو په منطق کښې به يې د دوه‌ځايي
نښلوونکي په توګه ورزياتوو۔ نو د $!A \lif
!B$ د بڼې پر !!{formula}s سربېره، د~$!A \cif !B$ د بڼې
!!{formula}s هم لرو۔ رسمي معنٰی پوهنه، د مودالي منطق د اړيکو د
معنٰی پوهنې په څېر، پر داسې مدلونو ولاړه ده چې !!{formula}s پکښې
په نړۍو کښې ارزول کېږي؛ او د $\mSat{M}{!A \cif !B}[w]$ د
تعريفوونکي تحقق شرط د ځينو (نورو) نړۍو~$w'$ لپاره د
$\mSat{M}{!A}[w']$ او $\mSat{M}{!B}[w']$ له مخې ورکول کېږي۔
کومې $w'$؟ د لومړني تصور له مخې، هغه يو يا څو چې~$w$ ته تر
ټولو نژدې وي او~$\mSat{M}{!A}[w']$ پکښې برقرار وي۔ دا غواړي
چې په مدل کښې د «نژدې‌والي» اړيکه هم شامله وي۔

لويس د واقعيت خلاف حالتونو د انځوريز ښودلو يوه ښوونکې طريقه
معرفي کړه۔ هره ممکنه نړۍ د يو داسې سټ په مرکز کښې ده چې يو
په بل کښې دننه کروي سيمې لري، او نورې نړۍ رانغاړي---مونږ دغه
کروي سيمې د هم‌مرکزو دايرو په توګه رسموو۔ د دوو کروي سيمو
ترمنځ نړۍ د مرکز نړۍ ته يو شان نژدې دي؛ په دننۍ کروي سيمه
کښې نړۍ نژدې، او په چاپېره کروي سيمه کښې نړۍ ليرې دي۔
\begin{center}
\begin{tikzpicture}[scale=.7]\small
  \spheresystem{5}
  \draw (0,0) node {$w$};
  \propositionintersect{0}{4}{40}{3.5}
  {\tikzset{proposition/.append={smooth,tension=1.4}}
  \proposition[shift={(1.6,-1)}]{90}{1}{30}{4}}
  \path (1.5,3) node[below] {$\formula{B}$};
  \path (3,0) node {$\formula{A}$};
\end{tikzpicture}
\end{center}
تر ټولو نژدې $!A$-نړۍ هغه نړۍ $w'$ دي چې~$!A$ پکښې تحقق
لري او د مرکز نړۍ~$w$ چاپېره په تر ټولو کوچنۍ کروي سيمه کښې
واقع دي (خړه سيمه)۔ د لومړني تصور له مخې، $!A \cif !B$
په~$w$ کښې هله تحقق لري چې $!B$ په ټولو تر ټولو نژدې
$!A$-نړۍو کښې رښتينے وي۔

\end{document}
